% GENERATED FILE. Edit canonical lessons, then run npm run generate:books.
% canonical-source-hash: fnv1a64:110cfbc55ac899f1
% canonical-lessons: PA-C177-noun-gender, PA-C177-adjective-position, PA-C177-adjective-agreement, PA-C177-bahut, PA-C177-noun-phrase

\chapter{A Good Boy and a Very Good Girl}
\label{ch:pa-noun-phrase-agreement-a1}
\hlchaptermodality{177}

\begin{chapteropening}
\textbf{By the end of this chapter:} \emph{I can build one masculine and one feminine adjective-noun phrase, intensify the quality, and keep the resulting noun phrase before the verb.}

Retrieve eh bahut changi kuri hai through listening, speaking, reading, and bounded Gurmukhi writing after building noun gender, adjective position, agreement, and degree one joint at a time.
\end{chapteropening}

\section[\texorpdfstring{\pa{ਮੁੰਡਾ} · \pa{ਕੁੜੀ} (munda · kuri)}{ਮੁੰਡਾ · ਕੁੜੀ (munda · kuri)}]{Two familiar nouns, now with their gender}

\label{lesson:PA-C177-noun-gender}



\begin{quote}
From the last chapter, say \textbf{\pa{ਕਿਤਾਬ} \pa{ਘਰ} \pa{ਵਿੱਚ} \pa{ਹੈ}} once. Now say the two familiar people words: \textbf{\pa{ਮੁੰਡਾ}} (\emph{muṇḍā}, boy), then \textbf{\pa{ਕੁੜੀ}} (\emph{kuṛī}, girl). Nothing about their meaning or spelling is new.
\end{quote}

\begin{grammarlens}[title={Grammar Lens: masculine and feminine nouns}]
Punjabi nouns have grammatical gender. For this first bounded pair:

\noindent
\begin{tabularx}{\linewidth}{@{}>{\raggedright\arraybackslash}X>{\raggedright\arraybackslash}X>{\raggedright\arraybackslash}X@{}}
\toprule
\textbf{noun} & \textbf{meaning} & \textbf{gender} \\
\midrule
\textbf{\pa{ਮੁੰਡਾ}} — \emph{muṇḍā} & boy & masculine \\
\textbf{\pa{ਕੁੜੀ}} — \emph{kuṛī} & girl & feminine \\
\bottomrule
\end{tabularx}

The endings help with this pair, but do not turn them into a universal guessing rule. Punjabi has masculine and feminine nouns whose endings do not advertise their gender. Learn the gender with the noun.
\end{grammarlens}

\subsection*{Your turn}
\begin{itemize}
\raggedright
  \item \emph{Hear:} \textbf{\pa{ਮੁੰਡਾ}}; YOU SAY: masculine
  \item \emph{Hear:} \textbf{\pa{ਕੁੜੀ}}; YOU SAY: feminine
  \item \emph{Read it:} \textbf{\pa{ਮੁੰਡਾ} · \pa{ਕੁੜੀ}}; YOU NAME: masculine · feminine
\end{itemize}

\begin{tcolorbox}[breakable,colback=teal!4,colframe=teal!35!black,title={Before you move on}]
Which familiar noun is masculine? (\textbf{\pa{ਮੁੰਡਾ}}.) Which is feminine? (\textbf{\pa{ਕੁੜੀ}}.) Keep the claim to this pair for now.

Source: University of Waterloo, \emph{PUN 1100 Course Package} (2018), pp. 54–55, \href{https://openlibrary-repo.ecampusontario.ca/jspui/bitstream/123456789/492/2/2018Pun1100CoursePackage.pdf}{Nouns and Gender}.
\end{tcolorbox}
\section[\texorpdfstring{\pa{ਚੰਗਾ} \pa{ਮੁੰਡਾ} (changa munda)}{ਚੰਗਾ ਮੁੰਡਾ (changa munda)}]{\pa{ਚੰਗਾ} \pa{ਮੁੰਡਾ} — put the quality before the noun}

\label{lesson:PA-C177-adjective-position}



\begin{quote}
First retrieve the last chapter's location frame once: \textbf{\pa{ਕਿਤਾਬ} \pa{ਘਰ} \pa{ਵਿੱਚ} \pa{ਹੈ}}. Then retrieve \textbf{\pa{ਚੰਗਾ}} (\emph{changā}, good) and \textbf{\pa{ਮੁੰਡਾ}} (\emph{muṇḍā}, boy). Name \textbf{\pa{ਮੁੰਡਾ}} as masculine. Keep the last two words separate for one breath.
\end{quote}

\begin{grammarlens}[title={Grammar Lens: one attributive position}]
Join them in this order:

\begin{quote}
\textbf{\pa{ਚੰਗਾ} \pa{ਮੁੰਡਾ}} — \emph{changā muṇḍā} — \textbf{good boy}
\end{quote}

In this noun phrase, \textbf{\pa{ਚੰਗਾ}} describes \textbf{\pa{ਮੁੰਡਾ}} and sits immediately before it: adjective first, noun second. This lesson owns that attributive pattern. It does not claim that an adjective can never appear after a noun in another kind of Punjabi sentence.
\end{grammarlens}

\subsection*{Your turn}
\begin{itemize}
\raggedright
  \item \emph{Build:} \textbf{\pa{ਚੰਗਾ}} + \textbf{\pa{ਮੁੰਡਾ}}
  \item \emph{Say it:} \textbf{\pa{ਚੰਗਾ} \pa{ਮੁੰਡਾ}}
  \item \emph{Point to:} adjective \textbf{\pa{ਚੰਗਾ}} first; noun \textbf{\pa{ਮੁੰਡਾ}} second
\end{itemize}

\begin{tcolorbox}[breakable,colback=teal!4,colframe=teal!35!black,title={Before you move on}]
Say \emph{good boy}. (\textbf{\pa{ਚੰਗਾ} \pa{ਮੁੰਡਾ}}.) Which word comes immediately before the noun? (\textbf{\pa{ਚੰਗਾ}}, the adjective.)

Source: Columbia University Punjabi Language Resource Center, \href{https://punjabi.lrc.columbia.edu/?page\_id=101}{Simple Sentence Structure}.
\end{tcolorbox}
\section[\texorpdfstring{\pa{ਚੰਗਾ} \pa{ਮੁੰਡਾ} · \pa{ਚੰਗੀ} \pa{ਕੁੜੀ} (changa munda …)}{ਚੰਗਾ ਮੁੰਡਾ · ਚੰਗੀ ਕੁੜੀ (changa munda …)}]{\pa{ਚੰਗਾ} \pa{ਮੁੰਡਾ} · \pa{ਚੰਗੀ} \pa{ਕੁੜੀ} — let the adjective agree}

\label{lesson:PA-C177-adjective-agreement}



\begin{quote}
Say \textbf{\pa{ਚੰਗਾ} \pa{ਮੁੰਡਾ}}. Then retrieve the familiar feminine noun \textbf{\pa{ਕੁੜੀ}}. Keep the adjective before the noun in both phrases.
\end{quote}

\begin{grammarlens}[title={Grammar Lens: change the adjective ending}]
\textbf{\pa{ਚੰਗਾ}} is a variable adjective. In this singular direct contrast, its ending agrees with the noun's gender:

\noindent
\begin{tabularx}{\linewidth}{@{}>{\raggedright\arraybackslash}X>{\raggedright\arraybackslash}X>{\raggedright\arraybackslash}X@{}}
\toprule
\textbf{noun gender} & \textbf{adjective + noun} & \textbf{meaning} \\
\midrule
masculine & \textbf{\pa{ਚੰਗਾ} \pa{ਮੁੰਡਾ}} — \emph{changā muṇḍā} & good boy \\
feminine & \textbf{\pa{ਚੰਗੀ} \pa{ਕੁੜੀ}} — \emph{changī kuṛī} & good girl \\
\bottomrule
\end{tabularx}

The noun changes from \textbf{\pa{ਮੁੰਡਾ}} to \textbf{\pa{ਕੁੜੀ}}; the adjective changes from \textbf{\pa{ਚੰਗਾ}} to \textbf{\pa{ਚੰਗੀ}}. This is one bounded singular direct contrast, not the whole Punjabi adjective system. Plural and oblique forms wait.
\end{grammarlens}

\subsection*{Your turn}
\begin{itemize}
\raggedright
  \item \emph{Hear:} \textbf{\pa{ਮੁੰਡਾ}}; YOU BUILD: \textbf{\pa{ਚੰਗਾ} \pa{ਮੁੰਡਾ}}
  \item \emph{Hear:} \textbf{\pa{ਕੁੜੀ}}; YOU BUILD: \textbf{\pa{ਚੰਗੀ} \pa{ਕੁੜੀ}}
  \item \emph{Switch:} \textbf{\pa{ਚੰਗਾ} \pa{ਮੁੰਡਾ} $\to$ \pa{ਚੰਗੀ} \pa{ਕੁੜੀ}}
\end{itemize}

\begin{tcolorbox}[breakable,colback=teal!4,colframe=teal!35!black,title={Before you move on}]
Say the masculine phrase, then the feminine phrase. Which piece changes to agree? (\textbf{\pa{ਚੰਗਾ} / \pa{ਚੰਗੀ}}, the adjective.)

Sources: University of Waterloo, \emph{PUN 1100 Course Package} (2018), pp. 71–74, \href{https://openlibrary-repo.ecampusontario.ca/jspui/bitstream/123456789/492/2/2018Pun1100CoursePackage.pdf}{Adjectives}; Columbia University Punjabi Language Resource Center, \href{https://punjabi.lrc.columbia.edu/?page\_id=101}{Simple Sentence Structure}.
\end{tcolorbox}
\section[\texorpdfstring{\pa{ਬਹੁਤ} (bahut)}{ਬਹੁਤ (bahut)}]{\pa{ਬਹੁਤ} — turn good into very good}

\label{lesson:PA-C177-bahut}



\begin{quote}
Say \textbf{\pa{ਚੰਗਾ} \pa{ਮੁੰਡਾ}}, then \textbf{\pa{ਚੰਗੀ} \pa{ਕੁੜੀ}}. Hear the one ending that changes and the adjective-before-noun order that stays.
\end{quote}

\subsection*{You'll want to know: \pa{ਬਹੁਤ}}
\begin{quote}
\textbf{\pa{ਬਹੁਤ}} — \emph{bahut} — \textbf{very}
\end{quote}

Put \textbf{\pa{ਬਹੁਤ}} before the adjective. It increases the degree; it does not take the adjective's agreement ending:

\begin{quote}
\textbf{\pa{ਬਹੁਤ} \pa{ਚੰਗਾ} \pa{ਮੁੰਡਾ}} — \emph{bahut changā muṇḍā} — very good boy
\end{quote}

\begin{quote}
\textbf{\pa{ਬਹੁਤ} \pa{ਚੰਗੀ} \pa{ਕੁੜੀ}} — \emph{bahut changī kuṛī} — very good girl
\end{quote}

The order is degree word + agreeing adjective + noun. Only \textbf{\pa{ਚੰਗਾ} / \pa{ਚੰਗੀ}} changes with this noun pair; \textbf{\pa{ਬਹੁਤ}} stays \textbf{\pa{ਬਹੁਤ}}.

\subsection*{Your turn}
\begin{itemize}
\raggedright
  \item \emph{Say it:} \textbf{\pa{ਚੰਗਾ} \pa{ਮੁੰਡਾ} $\to$ \pa{ਬਹੁਤ} \pa{ਚੰਗਾ} \pa{ਮੁੰਡਾ}}
  \item \emph{Say it:} \textbf{\pa{ਚੰਗੀ} \pa{ਕੁੜੀ} $\to$ \pa{ਬਹੁਤ} \pa{ਚੰਗੀ} \pa{ਕੁੜੀ}}
  \item \emph{Check:} \textbf{\pa{ਬਹੁਤ}} first; agreeing adjective second; noun last
\end{itemize}

\begin{tcolorbox}[breakable,colback=teal!4,colframe=teal!35!black,title={Before you move on}]
Make \textbf{\pa{ਚੰਗੀ} \pa{ਕੁੜੀ}} stronger. (\textbf{\pa{ਬਹੁਤ} \pa{ਚੰਗੀ} \pa{ਕੁੜੀ}}.) Which word adds \emph{very}? (\textbf{\pa{ਬਹੁਤ}}.) Which word still agrees? (\textbf{\pa{ਚੰਗੀ}}.)

Source: University of Waterloo, \emph{PUN 1100 Course Package} (2018), pp. 72–74, \href{https://openlibrary-repo.ecampusontario.ca/jspui/bitstream/123456789/492/2/2018Pun1100CoursePackage.pdf}{Adjective position and degree examples}.
\end{tcolorbox}
\section[\texorpdfstring{\pa{ਇਹ} \pa{ਬਹੁਤ} \pa{ਚੰਗੀ} \pa{ਕੁੜੀ} \pa{ਹੈ} (eh bahut changi …)}{ਇਹ ਬਹੁਤ ਚੰਗੀ ਕੁੜੀ ਹੈ (eh bahut changi …)}]{One phrase before the verb}

\label{lesson:PA-C177-noun-phrase}



\begin{quote}
Build \textbf{\pa{ਬਹੁਤ} \pa{ਚੰਗੀ} \pa{ਕੁੜੀ}} one piece at a time. Bring back \textbf{\pa{ਇਹ}} for \emph{this} and \textbf{\pa{ਹੈ}} for \emph{is}. Keep \textbf{\pa{ਹੈ}} waiting for the end.
\end{quote}

\begin{grammarlens}[title={Grammar Lens: the words work as one noun phrase}]
\begin{quote}
\textbf{\pa{ਇਹ} \pa{ਬਹੁਤ} \pa{ਚੰਗੀ} \pa{ਕੁੜੀ} \pa{ਹੈ।}}
\end{quote}

\begin{quote}
\emph{eh bahut changī kuṛī hai.}
\end{quote}

\begin{quote}
This is a very good girl.
\end{quote}

Before \textbf{\pa{ਹੈ}}, the words \textbf{\pa{ਬਹੁਤ} \pa{ਚੰਗੀ} \pa{ਕੁੜੀ}} travel together as a multiword noun phrase: degree word + agreeing adjective + noun. The noun is feminine, so the variable adjective is \textbf{\pa{ਚੰਗੀ}}. The verb remains last.

This proves only this one bounded singular direct pattern. It does not add noun plural, genitive, an article, or a rule for every Punjabi adjective.
\end{grammarlens}

\subsection*{Your turn}
\begin{itemize}
\raggedright
  \item \emph{Hear:} \emph{eh bahut changī kuṛī hai}; YOU POINT: \textbf{\pa{ਇਹ} \pa{ਬਹੁਤ} \pa{ਚੰਗੀ} \pa{ਕੁੜੀ} \pa{ਹੈ}}
  \item \emph{Say it:} \textbf{\pa{ਇਹ} \pa{ਬਹੁਤ} \pa{ਚੰਗੀ} \pa{ਕੁੜੀ} \pa{ਹੈ}}
  \item \emph{Read it:} \textbf{\pa{ਇਹ} \pa{ਬਹੁਤ} \pa{ਚੰਗੀ} \pa{ਕੁੜੀ} \pa{ਹੈ}}
  \item \emph{Write it:} \textbf{\pa{ਇਹ} \pa{ਬਹੁਤ} \pa{ਚੰਗੀ} \pa{ਕੁੜੀ} \pa{ਹੈ}} after covering the model
  \item \emph{Check:} \textbf{\pa{ਬਹੁਤ}} + \textbf{\pa{ਚੰਗੀ}} + \textbf{\pa{ਕੁੜੀ}} all come before \textbf{\pa{ਹੈ}}
\end{itemize}

\begin{tcolorbox}[breakable,colback=teal!4,colframe=teal!35!black,title={Before you move on}]
Rebuild the sentence from \emph{this}, \emph{very}, \emph{good}, \emph{girl}, \emph{is}. Name the multiword noun phrase before the verb. (\textbf{\pa{ਬਹੁਤ} \pa{ਚੰਗੀ} \pa{ਕੁੜੀ}}.)

This closes five named points in a partial project-defined inventory. It is not an A1 mock, not A1 readiness or attainment, and not external or human validation.

Sources: University of Waterloo, \emph{PUN 1100 Course Package} (2018), pp. 54–55, 71–74, \href{https://openlibrary-repo.ecampusontario.ca/jspui/bitstream/123456789/492/2/2018Pun1100CoursePackage.pdf}{Nouns and adjectives}; Columbia University Punjabi Language Resource Center, \href{https://punjabi.lrc.columbia.edu/?page\_id=101}{Simple Sentence Structure}.
\end{tcolorbox}
